\documentclass[12pt]{article}
\usepackage[utf8]{inputenc}
\usepackage[T1]{fontenc}
\usepackage[margin=1in]{geometry}
\usepackage{setspace}
\usepackage{xcolor}
\usepackage[colorlinks=true, urlcolor=blue!60!black, linkcolor=blue!60!black]{hyperref}
\setlength{\parindent}{0em}
\setlength{\parskip}{0.45em}
\renewcommand{\baselinestretch}{1.0}
\begin{document}
\begin{center}
{\Large \textit{Brandenburg v. Ohio}, 395 U.S. 444 (1969)\par}
\vspace{0.3em}
{\normalsize Case Brief --- One-Page Edition\par}
\end{center}

\textbf{Memory Jogger.} A Klan speech hinting at ``revengeance'' was protected, and the imminent-incitement test replaced \textit{Whitney} (\href{https://supreme.justia.com/cases/federal/us/395/444/}{Justia}).

\textbf{Facts.} At a filmed Klan rally on a farm, attended only by Klan members and a TV crew, Clarence Brandenburg said ``revengeance'' might be taken against the government and announced a march on Congress.

\textbf{Procedural History.} Convicted under Ohio's Criminal Syndicalism Statute for advocating violence as a means of political reform and sentenced to one to ten years; the Supreme Court of Ohio dismissed his appeal.

\textbf{Votes.} Reversed, 8--0, \emph{per curiam}; Black and Douglas concurring.

\textbf{Holding.} Advocacy of force may be punished only if directed to inciting imminent lawless action and likely to produce it; \textit{Whitney v.\ California} is overruled.

\textbf{Reasoning.} Ohio's law punished mere advocacy and assembly without distinguishing ``mere abstract teaching'' from preparing a group for violence. Black and Douglas rejected ``clear and present danger'' altogether.

\textbf{Examples.} \emph{Same-Side:} a pamphlet urging future revolution. \emph{Opposite-Side:} a bullhorn command to an armed crowd to storm a building now. \emph{Fence-Sitter:} a livestream call to raid a courthouse tomorrow.

\textbf{Critique.} The Court recast \textit{Dennis}'s balancing as an imminence test, left ``imminent'' and ``lawless action'' undefined, and adopted a live-crowd model that fits online incitement poorly.

\textbf{Key Quotation.} ``[T]he constitutional guarantees of free speech and free press do not permit a State to forbid or proscribe advocacy of the use of force or of law violation except where such advocacy is directed to inciting or producing imminent lawless action and is likely to incite or produce such action.'' --- \emph{Per curiam}, 395 U.S.\ at 447.

\end{document}
